\documentclass[10pt,a4paper]{amsart}
\usepackage[margin=19mm]{geometry}
\usepackage{amsmath,amssymb,mathtools}
\usepackage[scr=rsfs,cal=euler]{mathalfa}
\usepackage{xfrac}
\usepackage[unicode,hidelinks,bookmarks=false]{hyperref}
\newcommand{\ie}{i.e.\ }
\newcommand{\cf}{cf.\ }
\newcommand{\resp}{respectively\ }
\newcommand{\Subsection}{\subsection}
\providecommand{\nobreakdash}{\nobreak\hbox{-}}
\setlength{\emergencystretch}{2em}
\multlinegap0pt
\usepackage{amssymb}
\usepackage{tikz}
\usepackage{algorithm}\usepackage{algpseudocode}
\usepackage{enumerate}
\usepackage{xcolor}
\newtheorem{theorem}{Theorem}
\newtheorem{corollary}{Corollary}
\newcommand{\norm}[1]{\left\lVert#1\right\rVert}
\def\qsrank{\mathrm{qsrank}}
\newcommand{\upd}[1]{{\color{black}#1}} % updated text for the revision
\title{On the data-sparsity of the solution of Riccati equations with applications to feedback control}
\author{Stefano Massei,, Luca Saluzzi}
\date{Source: arXiv:2408.16569v2 (CC BY 4.0)}
\begin{document}
\maketitle

\noindent\emph{Source and licence.} On the data-sparsity of the solution of Riccati equations with applications to feedback control. Version arXiv:2408.16569v2 (CC BY 4.0), distributed by arXiv under the Creative Commons Attribution 4.0 International licence, http://creativecommons.org/licenses/by/4.0/, as recorded in the arXiv licence metadata for this exact version and shown on the abstract page https://arxiv.org/abs/2408.16569v2. Full licence notice: \texttt{licenses/arxiv-2408.16569-CC-BY-4.0.txt}.\par\smallskip

\noindent\textit{Editorial scope.} Counted unit: the article's corollary cor:sing-decay3 (source0.tex lines 596--608). The article's complete original proof of it is delivered with it (source0.tex lines 609--648, one \texttt{proof} environment of 40 lines). No mathematics is written, repaired or re-derived: the statement and the proof are the article's own lines, in the article's own notation, and no retained line is substituted or re-flowed.  Retained uncounted as the source-local context the proof needs: lines 119--121; lines 386--422; lines 529--538.  Nothing was omitted from any retained span, so the package carries no omission record.  No retained line contains a citation, so the wrapper carries no bibliography and no bibliography entry.  No retained line contains a figure environment, an image-inclusion command or drawing code, so no image is delivered and the package declares no asset dependency.  Editorial formatting only, in addition: an AMS \texttt{amsart} preamble that restates the article's own declarations and macros used by the retained lines, namely the environments \texttt{theorem}, \texttt{corollary} on separate counters, together with the standard packages those lines need.  The retained environments print in this document as Theorem 1, Corollary 1 in order of appearance, whereas the article prints the same items with the numbering of its own full document.  The complete native source of the article as distributed in the arXiv e-print for this exact version is bundled unchanged as \texttt{sources/164453/source0.tex}.
\begin{equation}\label{eq:care}
\mathcal{R}(X) := A^\top X+XA- XFX+Q=0,	
\end{equation}

\begin{theorem}\label{thm:sing-decay}
	Let $H\in\mathbb R^{2n\times 2n}$ be of the form
	$$
	H=\begin{bmatrix}
		A&-F\\
		-Q&-A^\top
	\end{bmatrix}
	$$
	where $A\in\mathbb R^{n\times n}$ is  quasiseparable of order $r_a$, $\mathcal W(A)\subset\mathbb C^{-}$, $F,Q\in\mathbb R^{n\times n}$ symmetric positive semidefinite  and quasiseparable of order $r_f$ and $r_q$, respectively, and let $X\in\mathbb R^{n\times n}$ be the symmetric positive semidefinite stabilizing solution of the continuous-time Riccati equation 
	$$
	A^\top X+XA-XFX+Q=0.
	$$
	Moreover, let $H^{(j)},H_{(j)}\in\mathbb R^{2j\times 2j}$ be defined as
	$$
	H^{(j)}=\begin{bmatrix}
		A^{(j)}&-F^{(j)}\\
		-Q^{(j)}&-(A^{(j)})^\top
	\end{bmatrix},\qquad H_{(j)}=\begin{bmatrix}
		A_{(j)}&-F_{(j)}\\
		-Q_{(j)}&-(A_{(j)})^\top
	\end{bmatrix},
	$$
	where $A^{(j)},F^{(j)},Q^{(j)}$  are the principal submatrices of $A,F,Q$ corresponding to the first $j$ indices of rows and columns, and $A_{(j)},F_{(j)},Q_{(j)}$ are those corresponding to the last $n-j$, for $j=1,\dots, n$. Finally, indicate with $X^{(j)}$ and $X_{(j)}$ the symmetric positive semidefinite stabilizing solutions of the Riccati equations associated with $H^{(j)}$ and $H_{(j)}$, respectively. 
	If $E\subset \mathbb C^-$ is such that
	$$
	\mathrm{convex\ hull}\left(\bigcup_{j=1,\dots, n}\mathcal W(A^{(j)} - F^{(j)}X^{(j)}) \quad \cup  \bigcup_{j=1,\dots,n}\mathcal W(A_{(j)} - F_{(j)}X_{(j)})\right)\subset E,
	$$
	then the singular values $\sigma_i(M)$ of any $M\in\mathrm{Off}(X)$ satisfy
	%$$
	%\frac{\sigma_{hk+1}}{\norm{ X}_2}\le (1+\sqrt 2)^2 \min_{r(z)\in\mathcal R_{h,h}}\frac{\max_{z\in E} |r(z)|}{\min_{z\in -E} |r(z)|}, \quad h=1,2,\dots 
	%$$
	%where $k=4b_a+2b_f+2b_q$, and $\mathcal R_{h,h}$ indicates the set of rational functions with both numerator and denominator of degree at most $h$.
	$$
	\frac{\sigma_{ht+1}(M)}{\norm{X}_2}\le (1+\sqrt 2)^2 Z_h(E,-E), \quad h=1,2,\dots, 
	$$
	where $t=4r_a+2r_f+2r_q$.
\end{theorem}

\begin{corollary}\label{cor:sing-decay2}
	Under the same assumptions of Theorem~\ref{thm:sing-decay}, if $F=I$, then the singular values of any $M\in\mathrm{Off}(X)$ satisfy
	$$
	\frac{\sigma_{ht+1}(M)}{\norm{X}_2}\le (1+\sqrt{2})^2
	Z_h(E,-E)
	%\min_{r(z)\in\mathcal R_{h,h}}\frac{\max_{z\in G} |r(z)|}{\min_{z\in -G} |r(z)|}, 
	\quad h=1,2,\dots,
	$$
	where $t=4r_a+2r_q$, $E=\mathcal W(A) + \left[-\tau -\sqrt{\tau^2+\norm{Q}_2},\ 0\right]$, and $\tau:= |\min\{\mathcal W(A)\cap \mathbb R\}|$.
\end{corollary}

\begin{corollary}\label{cor:sing-decay3}
	Under the same assumptions of Theorem~\ref{thm:sing-decay}, if $F$ is symmetric positive definite, $LL^\top=F$ denotes its Cholesky factorization, and $\mathcal W(L^{-1}AL)\subset \mathbb C^-$,  then the singular values of  any $M\in\mathrm{Off}(X)$   satisfy
	$$
	\frac{\sigma_{ht+r_f+1}(M)}{\norm{X}_2}\le \kappa(F)(1+\sqrt{2})^2 Z_h(E,-E),
	%\min_{r(z)\in\mathcal R_{h,h}}\frac{\max_{z\in G} |r(z)|}{\min_{z\in -G} |r(z)|},
	\quad h=1,2,\dots,
	$$
	where $$t=4r_a+6r_f+2r_q,$$
	 $$
	 E=\mathcal W(L^{-1}AL) + \left[-\tau -\sqrt{\tau^2+\norm{F}_2\norm{Q}_2},\ 0\right],
	 $$
	   $$\tau:= |\min\{\mathcal W(L^{-1}AL)\cap \mathbb R\}|.$$   
\end{corollary}

\begin{proof}
	Multiplying \eqref{eq:care} from the left by $L^\top$ and from the right by $L$ we get
	\begin{equation}\label{eq:riccati-cholesky}
		\widetilde A^\top Y + Y\widetilde A- Y^2+ \widetilde Q=0,
	\end{equation}
	with $\widetilde A:=L^{-1}AL, \widetilde Q =L\upd{^\top}QL$, and $Y:=L^\top XL$. Note that $L$ is lower triangular and quasiseparable of order $r_f$  so that also $L^{-1}$ is lower triangular and quasiseparable with the same order. In view of the \upd{subadditivity} properties of the quasiseparable rank, we have that $\qsrank(\widetilde A)\le r_a+r_f$, and  $\qsrank(\widetilde Q)\le r_f+r_q$.
	By applying Corollary~\ref{cor:sing-decay2} to \eqref{eq:riccati-cholesky} we get 
	\begin{equation}\label{eq:Ybound}
		\frac{\sigma_{ht+1}(M_Y)}{\norm{Y}_2}\le (1+\sqrt{2})^2 Z_h(E,-E),
		%\min_{r(z)\in\mathcal R_{h,h}}\frac{\max_{z\in G} |r(z)|}{\min_{z\in -G} |r(z)|}
	\end{equation}
	where $M_Y$ indicates 
	any element of $\mathrm{Off}(Y)$. Observe that if we look at the $2\times 2$ block partitioning
	$$
	Y=\begin{bmatrix}
		Y_{11} & Y_{12}\\
		Y_{21}&Y_{22}
	\end{bmatrix}, \qquad X=\begin{bmatrix}
		X_{11} & X_{12}\\
		X_{21}&X_{22}
	\end{bmatrix}\qquad L^{-1}=\begin{bmatrix}
		\widetilde L_{11} & 0\\
		\widetilde L_{21}&\widetilde L_{22}
	\end{bmatrix}
	$$
	where the diagonal blocks are square, and all the matrices are partitioned in the same way, then we have that the offdiagonal blocks of $X=L^{-\upd{\top}}YL^{-1}$ are given by 
	$$
	X_{21}= \widetilde L_{22}^\top Y_{21}\widetilde L_{11} +\widetilde L_{22}^\top Y_{22}\widetilde L_{21}, \qquad  X_{12}= \widetilde L_{11}^\top Y_{12}\widetilde L_{22} +\widetilde L_{21}^\top Y_{22}\widetilde L_{22}.
	$$
	Let us show that the singular values of $X_{21}$ satisfy the claim. Since the rank of $\widetilde L_{22}^\top Y_{22}\widetilde L_{21}$ is at most $r_f$ and the singular values of $Y_{21}$ satisfy \eqref{eq:Ybound}, then we have
	\begin{align*}
		\sigma_{ht+r_f+1}(X_{21})\le \sigma_{hk+1}(\widetilde L_{22}^\top Y_{21}\widetilde L_{11})&\le \norm{F^{-1}}_2 \sigma_{ht+1}(Y_{21})\\ &\le \norm{F^{-1}}_2\norm{Y}_2 (1+\sqrt{2})^2 Z_h(E,-E)
		%\min_{r(z)\in\mathcal R_{h,h}}\frac{\max_{z\in G} |r(z)|}{\min_{z\in -G} |r(z)|} 
		\\ & \le \norm{F^{-1}}_2 \norm{F}_2\norm{X}_2 (1+\sqrt{2})^2
		Z_h(E,-E),
		%\min_{r(z)\in\mathcal R_{h,h}}\frac{\max_{z\in G} |r(z)|}{\min_{z\in -G} |r(z)|},
	\end{align*}
	where we have used the relations: $\norm{\widetilde L_{jj}}_2\le \norm{L^{-1}}_2$, $\norm{L^{-1}}_2\norm{L^{-\upd{\top}}}_2= \norm{F^{-1}}_2$, and $\norm{L}_2\norm{L^\top}_2= \norm{F}_2$.
	As the matrix $X$ is symmetric, there is no need to repeat the argument for $X_{12}$,  and this concludes the proof. 
\end{proof}

\end{document}
